\chapter{总结与展望}
\label{chap:conclusion}

本文以百度Apollo 11.0开源自动驾驶平台为基础，针对多障碍物密集场景下路径边界构建保守或失效的问题，提出多障碍物交互运动学统一规划算法，即MIKU（Multi-obstacle Interaction-density Kinematic-prediction Unified Planner）。源码分析显示，历史版本中的\texttt{PathBoundsDecider}在Apollo 11.0中已主要下沉为\texttt{PathBoundsDeciderUtil}工具类，其边界构建逻辑表现出三个特征，即逐障碍物独立决策、统一安全裕度，以及动态障碍物在路径阶段被滤除。笔者据此从安全距离、决策粒度与时空协调三个方向切入改进。

算法层面，本文在SL层建立融合TTC、横向重叠率、相对速度、障碍物类型与交互密度的威胁度模型，以差异化安全裕度$\delta_i$释放低威胁障碍物占用的通行空间。针对组内避让方向选择，本文将问题建模为最大通行带宽度优化，证明了最优分割在横向排序下的连续性以及最大间隙策略在该形式化下的最优性，将组内搜索从$O(2^k)$候选压缩为$k+1$个候选分割，精确求解复杂度为$O(k \log k)$。在ST层，本文利用到达时间函数$\tau(s)$与预测轨迹将动态障碍物纳入路径边界构建，并将时变通行带的时间有效性约束投影到速度QP中，与\texttt{SpeedBoundsDecider}的ST边界融合。MIKU总复杂度为扫描线分组的$O(n\log n)$叠加逐离散点边界构建的$O(K)$，合计$O(n\log n + K)$，其中$K$为路径或时间离散化规模；同场景下Baseline的$O(nK)$在$n$较大时劣于本文方法。

工程与实验层面，本文已在Apollo Planning C++代码中完成增量接入，通过\texttt{LaneFollowPath}新增MIKU路径边界分支，并以\texttt{STDrivableBoundary}形式向\texttt{PiecewiseJerkSpeedOptimizer}传递速度约束。Python复现的\MPressureScenarioTotal{}个压力场景与\MComparableScenarioTotal{}个可比场景显示，在压力场景中Baseline与MIKU的通过率分别为\MPressureBaselinePassCount{}/\MPressureScenarioTotal{}与\MPressureMikuPassCount{}/\MPressureScenarioTotal{}；在可比场景中两者均完成通行，MIKU在动态场景下降低了纵向加减速峰值，C1--C4平均QP求解总耗时由\MComparableBaselineQpTotal\,ms降至\MComparableMikuQpTotal\,ms。压力场景中Baseline进入FallbackPath降级并使QP提前中止，其耗时记录不构成效率比较依据。Dreamview在环展示了C++接入后的规划行为，实车可用性仍需在更完整的感知、预测与控制链路中验证。

后续仍有若干可检验的研究方向。第一，引入交互式预测，将自车候选行为反馈至周车预测过程，以降低开环预测在密集交互场景中的偏差。第二，将当前单次前向的到达时间估计扩展为路径与速度的迭代求解，并对迭代收敛性以及单帧计算预算给出明确界。第三，对纵向链式重叠或仅部分重叠的障碍物组进行更细粒度的截面求解，量化空间利用率与计算开销之间的权衡曲线。第四，当多个候选间隙宽度接近时引入学习辅助评分，补入变道代价、道路拓扑与驾驶习惯等因素。第五，在真实车辆平台上评估感知噪声、预测误差与规划-控制延迟对时变通行带的影响，并给出可复现的降级策略与失效安全方案。

本文验证限于Python复现与Dreamview在环阶段，未覆盖多车交互、恶劣天气与长时程漂移等工况，上述局限应在后续实车实验中进一步暴露并修正。
